\chapter{Double Complexes and Tensor Calculations}
\label{ch:double-complexes}

A double complex has two filtrations on its total complex. The two spectral
sequences compute the same total cohomology under a common finiteness
hypothesis, but the two filtrations on that cohomology need not coincide.
We use cohomological indices for the general constructions and switch to
homological indices for the tensor calculations. The sign and indexing
conventions below specify how to pass between the two.

\section{Totalization and its filtrations}

\begin{proposition}[Totalization and interchange of the axes]
\label{res:B01}
Let \(C^{p,q}\), \(p,q\in\mathbb Z\), be objects of an additive category
with differentials \(d_h\colon C^{p,q}\to C^{p+1,q}\) and
\(d_v\colon C^{p,q}\to C^{p,q+1}\). Suppose that
\(d_h^2=d_v^2=0\) and that the horizontal and vertical squares commute.
On \(C^{p,q}\) define
\[
    D=d_h+(-1)^p d_v .
\]
This gives a differential on the direct-sum totalization when the requisite
diagonal coproducts exist, and on the product totalization when the
diagonal products exist. Multiplying \(d_v\) by \((-1)^p\) changes to the
anticommuting convention and leaves \(D\) unchanged. The transposed double
complex \(C^{\mathrm t}\) has \((C^{\mathrm t})^{q,p}=C^{p,q}\) with its
horizontal and vertical differentials interchanged. Multiplication by
\((-1)^{pq}\) from \(C^{p,q}\) to \((C^{\mathrm t})^{q,p}\) induces a
natural isomorphism of total complexes in either totalization.
\end{proposition}

\begin{proof}
Write \(a_{p,q}\) for the injection of a summand and \(b_{p,q}\) for the
projection to a factor, with the total degree specified by \(p+q\).
For the direct sum, the outgoing differential of each summand is the
finite sum
\[
  D a_{p,q}=a_{p+1,q}d_h+(-1)^p a_{p,q+1}d_v.
\]
The coproduct property gives a unique map with these formulas. For
the product, the incoming map to each output factor is
\[
  b_{p,q}D=d_h b_{p-1,q}+(-1)^p d_v b_{p,q-1}.
\]
Only two input coordinates occur, so the product property again gives
a map. Composing twice and collecting the components gives the two
square terms \(d_h^2,d_v^2\) and the mixed terms
\[
     (-1)^{p+1}d_vd_h+(-1)^p d_hd_v
\]
on input bidegree \((p,q)\). They all vanish. The same finite component
calculation applies to the product by projecting to an output factor.
With \(\widetilde d_v=(-1)^p d_v\), the mixed terms express
\(d_h\widetilde d_v+\widetilde d_vd_h=0\). Rescaling by \((-1)^p\)
again recovers the commuting convention.

On the transposed totalization the differential, in the original
coordinates, is \(d_v+(-1)^q d_h\). If \(T\) multiplies the input
\(C^{p,q}\) by \((-1)^{pq}\), the coefficients of \(d_h\) and \(d_v\)
in \(TD\) are respectively
\((-1)^{(p+1)q}\) and \((-1)^p(-1)^{p(q+1)}\).
These agree with the coefficients of \(d_h\) and \(d_v\) in
\(D^{\mathrm t}T\). The formula defines an inverse by the same
multipliers and reindexing in the opposite direction. These formulas
commute with maps preserving both differential directions.
\end{proof}

The signs specialize to the tensor total complex: when both differentials
lower degree, the same formula \(D=d_h+(-1)^p d_v\) works, since replacing
\(p+1\) by \(p-1\) does not change the parity calculation.
See \cite[Sections 5.1 and 5.6]{weibel} and
\cite[Chapter 12, Sections 18 and 25]{stacks-homology}.

\begin{lemma}[Finite diagonals and the column and row filtrations]
\label{res:B02}
Let \(C\) be a double cochain complex in an abelian category and suppose
that for each \(n\) the set
\[
    I_n=\{(p,q)\colon p+q=n,\ C^{p,q}\ne 0\}
\]
is finite. The canonical map from the diagonal direct sum to the diagonal
product is an isomorphism of complexes. Write \(K=\operatorname{Tot}C\)
for their common value. There are two natural decreasing filtrations by
subcomplexes
\begin{align*}
    F_{\mathrm c}^{a} K^n
       &=\bigoplus_{\substack{p+q=n\\p\ge a}} C^{p,q},\\
    F_{\mathrm r}^{a} K^n
       &=\bigoplus_{\substack{p+q=n\\q\ge a}} C^{p,q}.
\end{align*}
Each filtration is finite in every degree. In particular, these statements
apply when \(C^{p,q}=0\) unless \(p\ge p_0\) and \(q\ge q_0\) for fixed
integers \(p_0,q_0\).
\end{lemma}

\begin{proof}
The nonzero factors are indexed by \(I_n\), and the remaining zero factors
may be removed by the universal properties of products and coproducts.
The diagonal comparison is the biproduct isomorphism whose \((i,j)\)-entry
is the identity for \(i=j\) and zero otherwise. The component formulas in
Proposition~\ref{res:B01} show that these comparisons intertwine the
differentials. In the column filtration, \(d_h\) increases the column
and \(d_v\) preserves it, so both carry the indicated summands into the
corresponding subobject in degree \(n+1\). In the row filtration, \(d_v\)
increases the row and \(d_h\) preserves it. Each resulting inclusion is
the split inclusion of a finite sub-biproduct in a given degree; the
splittings need not be maps of complexes.

For a nonzero diagonal, a threshold at most the minimum relevant index
gives the whole degree, and a threshold greater than the maximum gives
zero. The zero diagonal satisfies both bounds. For the translated
first quadrant, a nonzero summand in degree \(n\) must have
\(p_0\le p\le n-q_0\). This gives the required finiteness and allows
the bounds to depend on \(n\).
\end{proof}

Outside the finite-diagonal setting the two totalizations can have
different cohomology. For example, take \(C^{p,-p}=k\) for every integer
\(p\), with \(k\) a nonzero field, and all other terms and differentials
zero. The diagonal comparison on \(H^0\) is the inclusion of finite
support sequences in all sequences. The convergence statements below
use exactly the finite-diagonal hypothesis unless another bound is
explicitly given.

\section{The two spectral sequences}

Use \(H_v^q(C^{p,*})\) for the cohomology of the vertical complex in
column \(p\). The horizontal maps are maps of these unsigned vertical
complexes and induce a differential on \(p\mapsto H_v^q(C^{p,*})\).
The notation \(H_h^p H_v^q(C)\) refers to its degree \(p\) cohomology.
The notation for the opposite order is defined by transposing the axes.

\begin{proposition}[The first pages and their signs]
\label{res:B03}
Suppose \(C\) satisfies Lemma~\ref{res:B02}. For the column filtration
on \(\operatorname{Tot}C\) the first pages, in cohomological convention,
are
\begin{align*}
    E_{0,\mathrm c}^{p,q}&=C^{p,q},
        &d_{0,\mathrm c}^{p,q}&=(-1)^p d_v,\\
    E_{1,\mathrm c}^{p,q}&=H_v^q(C^{p,*}),
        &d_{1,\mathrm c}^{p,q}&=H_v^q(d_h),\\
    E_{2,\mathrm c}^{p,q}&=H_h^p H_v^q(C).
\end{align*}
For the row filtration use \(p\) for filtration degree, so that it is
the original row index. Through the signed interchange of
Proposition~\ref{res:B01} the row pages identify with the column pages
of \(C^{\mathrm t}\). In particular
\[
    E_{1,\mathrm r}^{p,q}\simeq H_h^q(C^{*,p}),
    \qquad
    E_{2,\mathrm r}^{p,q}\simeq H_v^p H_h^q(C).
\]
These identifications are natural for maps of double complexes.
\end{proposition}

\begin{proof}
In total degree \(p+q\) the quotient of the column terms at thresholds
\(p,p+1\) is the biproduct factor \(C^{p,q}\). The horizontal term of \(D\)
lands in the subobject being quotiented out, so the associated graded
differential is \((-1)^p d_v\). Multiplication of a differential by
\(-1\) leaves its kernel and image unchanged, so there is a canonical
identification of the first page with the unsigned vertical cohomology.
For a representative in column \(p\) that is vertically closed, its
total differential is \(d_h\) of the representative, now in column
\(p+1\). The commutation of the unsigned differentials implies that
\(d_h\) takes vertical cycles to cycles and vertical boundaries to
boundaries.

The same argument works in an abelian category without element notation.
The inclusion of the kernel of \(d_v\) in \(C^{p,q}\) has composite
with \(D\) factoring through the column \(p+1\) summand, and the
induced map to the quotient of cycles by boundaries is the map induced
by \(d_h\). This is the first filtered differential of
Theorem~\ref{res:F04}. Taking its cohomology gives the second page.

The interchange isomorphism takes \(F_{\mathrm r}\) for \(C\) to
\(F_{\mathrm c}\) for \(C^{\mathrm t}\). Equivalently, in the original
coordinates \(C^{q,p}\) on the raw row page has \(d_0=d_h\), and the raw
\(d_1\) is induced by \((-1)^q d_v\). The signed identification on
\(E_1^{p,q}\) multiplies by \((-1)^{pq}\). Increasing \(p\) by one changes
this multiplier by \((-1)^q\), which accounts for the raw sign.
Thus the two descriptions give isomorphic spectral sequences with the
same convention on later pages. The constructions use natural
sub-biproduct inclusions, kernels, and images.
\end{proof}

\begin{theorem}[Convergence, comparison, and assembly]
\label{res:B04}
Suppose \(C\) has finite diagonals. Both spectral sequences of
Proposition~\ref{res:B03} converge to \(H^*(\operatorname{Tot}C)\).
More precisely, each gives, in degree \(n\), the finite decreasing
filtration by images of
\[
  H^n(F^a\operatorname{Tot}C)\longrightarrow H^n(\operatorname{Tot}C),
\]
and \(E_\infty^{p,n-p}\) is naturally isomorphic to its \(p\)-th graded
piece. If a map \(C\to C'\) between two such double complexes is a
quasi-isomorphism on each column, or on each row, then the map on
totalizations is a quasi-isomorphism. In particular, exactness of every
column, or of every row, implies acyclicity of the total complex.

Suppose moreover \(C^{p,q}=0\) for \(q<0\), and a map of horizontal
complexes \(L^*\to C^{*,0}\) is annihilated by \(d_v\) and induces
\(L^p\simeq H_v^0(C^{p,*})\) for every \(p\). If the other vertical
cohomology vanishes, the induced map \(L^*\to\operatorname{Tot}C\) is
a quasi-isomorphism. The transposed assertion holds for an augmentation
into \(C^{0,*}\) when \(C^{p,q}=0\) for \(p<0\).
\end{theorem}

\begin{proof}
By Lemma~\ref{res:B02}, each filtration satisfies the degreewise finite
hypothesis of Theorem~\ref{res:S01}. Its conclusion supplies the
stabilized page and graded isomorphisms, including their compatibility
with the induced map on cohomology. A columnwise quasi-isomorphism is an
isomorphism on the \(E_{1,\mathrm c}\) page by
Proposition~\ref{res:B03}. Hence it is an isomorphism on the stabilized
page. In each total degree both induced filtrations are finite.
For a morphism between such filtered objects, descend through
\[
    0\longrightarrow F^{a+1}H^n
      \longrightarrow F^aH^n
      \longrightarrow \operatorname{gr}_F^a H^n
      \longrightarrow 0 .
\]
Starting at a large common zero threshold, an isomorphism on the first
and third terms gives one on the middle term by the short five lemma.
At a common lower threshold this gives the map on the whole cohomology
object. This is also the finite comparison of
Corollary~\ref{res:S03}. The row result follows by signed interchange,
and comparison with the zero double complex gives the acyclicity claim.

For the augmentation, view \(L^p\) as the double complex concentrated
in bidegrees \((p,0)\). The hypothesis about \(d_v\) makes the
augmentation a map of double complexes. It is a quasi-isomorphism on
each column: in degree zero it induces the given isomorphism, and in
the other degrees both cohomology objects vanish. Its totalization is
the stated map from \(L^*\). For the row case transpose and use
Proposition~\ref{res:B01}; on the axis the interchange multiplier is
\((-1)^0=1\).
\end{proof}

The same argument applies to chain complexes with all indices
negated. Suppose a first-quadrant chain double complex has a vertical
augmentation commuting with the horizontal maps, exact in positive
vertical degrees and inducing an isomorphism on vertical \(H_0\).
The augmentation annihilates boundaries from vertical degree one
and defines a map to a double complex concentrated in vertical
degree zero. It is a quasi-isomorphism on each column, so the induced
map from the total complex is a quasi-isomorphism as well.
An analogous statement holds
for horizontal resolutions. The module assembly and augmented
comparisons in \cite[Chapter 10]{p103} specialize this argument. Compare
\cite[Lemmas 12.25.1, 12.25.3, and 12.25.4]{stacks-homology} and
\cite[Exercise 5.6.4]{weibel}.

\section{Tensor complexes and change of rings}

For a chain double complex \(C_{p,q}\), put \(C^{i,j}=C_{-i,-j}\).
Descending column threshold \(i\ge -a\) then means original chain
column \(p\le a\), which defines the usual increasing filtration
\(F_a\) of the sum totalization. Replacing \((i,j)\) by \((-p,-q)\)
sends the cohomological page differential of bidegree
\((r,1-r)\) to a chain page differential of bidegree \((-r,r-1)\).
Thus the same proofs compute successive quotients
\(F_pH_{p+q}/F_{p-1}H_{p+q}\). For the row filtration we continue to put
its filtration index first, so that the displayed first index is the
original vertical degree in the chain double complex.

In this section rings have identity, and module actions preserve it.
We write the right \(R\)-module first in \(\otimes_R\) and in
\(\operatorname{Tor}^R_n\), and the left \(R\)-module second. Unless an
additional module action is specified, these constructions take values
in abelian groups. We use homological grading, with resolutions in
degrees \(0,1,\ldots\). The differential on the tensor double complex
of right and left complexes \(P,Q\) is
\[
   \partial(x\otimes y)=\partial_P x\otimes y
            +(-1)^i x\otimes\partial_Q y,\qquad x\in P_i.
\]
For a flat module, tensoring in the other variable is exact. In
particular, it carries the cycle and boundary short exact sequences to
short exact sequences and gives the canonical maps identifying
homology after tensoring with the tensor products of the homology
groups. This also gives the comparisons on the first pages used below.

\begin{theorem}[Balance of Tor]
\label{res:B05}
For a right \(R\)-module \(A\) and a left \(R\)-module \(M\), let
\(P_\bullet\to A\) and \(Q_\bullet\to M\) be projective resolutions in
their respective module categories. There is a natural comparison,
independent of the choices of projective resolutions,
\[
        H_n(P_\bullet\otimes_R M)
        \;\simeq\; H_n(A\otimes_R Q_\bullet).
\]
Both groups compute \(\operatorname{Tor}_n^R(A,M)\).
\end{theorem}

\begin{proof}
The double complex has terms \(P_i\otimes_R Q_j\) with the two commuting
differentials. Each \(P_i\) is a flat right module, so the augmented
column is exact after tensoring and computes \(P_i\otimes_R M\) in
degree zero. Thus the map from the total complex induced by
\(Q_\bullet\to M[0]\) is a quasi-isomorphism by
Theorem~\ref{res:B04}. Likewise each \(Q_j\) is a flat left module, so
augmentation by \(P_\bullet\to A[0]\) is a quasi-isomorphism from the
same total complex to \(A\otimes_R Q_\bullet\). The map from the first
homology group to the second is the inverse of the first total
comparison followed by the second. Both augmentations are chain maps:
they vanish away from the corresponding degree-zero row or column,
commute along that row or column, and annihilate boundaries entering
it. In the horizontal degree-zero quotient the tensor differential in
the remaining direction has sign \((-1)^0=1\).

For a morphism of module pairs, lift the two coefficient maps to maps
of projective resolutions. The comparison theorem for projective
resolutions gives lifts unique up to chain homotopy
\cite[Chapter 10]{p103}. If two lifts in the right variable differ
by \(\partial h+h\partial\), their maps on total complexes differ by
the homotopy \(x\otimes y\mapsto h(x)\otimes y\). Indeed the vertical
terms in its homotopy identity cancel since \(h\) increases \(i\) by one.
If two lifts in the left variable differ by
\(\partial h+h\partial\), use
\(x\otimes y\mapsto(-1)^i x\otimes h(y)\);
now the horizontal terms cancel. This proves independence on
homology. Both augmentations commute with the lifted pair of maps.
Applying the argument to lifts of identity maps compares two choices
of resolutions by inverse maps on homology. These identities prove
naturality and independence for the displayed comparison. The two
augmentations determine the comparison on homology. See
\cite[Chapter 11]{p103} for the module resolution models and
\cite[Application 5.6.3]{weibel} for the spectral-sequence argument.
\end{proof}

\begin{theorem}[The K\"unneth spectral sequence]
\label{res:B06}
Let \(P_\bullet\) be a chain complex of flat right \(R\)-modules with
\(P_i=0\) for \(i<L\) for some integer \(L\). For a left \(R\)-module \(M\)
there is a natural homological spectral sequence
\[
    E^2_{p,q}=\operatorname{Tor}^R_p(H_q(P),M)
       \quad\Longrightarrow\quad H_{p+q}(P_\bullet\otimes_R M).
\]
Its support lies in \(p\ge0,\ q\ge L\), with a finite target filtration
in each degree and differential \(d_r\) of bidegree \((-r,r-1)\).
\end{theorem}

\begin{proof}
Choose a projective left resolution \(Q_\bullet\to M\).
Put \(D_{i,j}=P_i\otimes_R Q_j\). In total degree \(n\), possible
nonzero terms have \(L\le i\le n,\ j=n-i\ge0\), so the
finite-diagonal comparison applies even when \(L<0\).
Flatness of \(P_i\) makes the vertical augmentation a
quasi-isomorphism on each column, giving the natural target
identification
\[
     H_n(\operatorname{Tot} D)\simeq H_n(P_\bullet\otimes_R M).
\]
For the other filtration, put \(p=j\) and \(q=i\).
Flatness of \(Q_p\) gives horizontal homology
\[
      E^1_{p,q}\simeq H_q(P)\otimes_R Q_p .
\]
Under the signed interchange that orders the \(Q\) index first, the
induced differential is \(1\otimes\partial_Q\). Equivalently, in the
untransposed coordinates it has the constant factor \((-1)^q\) in
each fixed \(q\). The cohomological sign calculation of
Proposition~\ref{res:B03} with the indices negated gives the isomorphism
of the two descriptions. Taking homology in \(p\) and using
Theorem~\ref{res:B05} gives the displayed \(E^2\). The finite-diagonal
bound and Theorem~\ref{res:B04} supply both convergence and the finite
filtration on the target.

A chain map of right complexes gives a map of the double complexes.
For a coefficient map in \(M\), use a lifted map of the projective
left resolutions. Two lifts differ by a homotopy which, after applying
\(H_q(P)\otimes_R -\), is an ordinary homotopy of the first-page
complexes in the \(p\)-direction. Hence they induce the same map on
\(E^2\), and therefore on all later pages since those maps are induced
on homology of pages. The augmentation comparisons commute with the
maps on the two targets. Taking lifts between two choices for the same
\(M\) yields natural comparisons from \(E^2\) onward and on the
filtered target. The resulting spectral sequence is the one of
\cite[Theorem 5.6.4]{weibel}.
\end{proof}

\begin{corollary}[The two-column case]
\label{res:B08}
In Theorem~\ref{res:B06}, suppose every \(H_q(P)\) admits a flat
resolution of length at most one as a right \(R\)-module. Then for
every \(n\) the target filtration gives a natural short exact sequence
\[
\begin{split}
 0 \longrightarrow H_n(P)\otimes_R M
    \longrightarrow H_n(P_\bullet\otimes_R M)
    \longrightarrow \operatorname{Tor}_1^R(H_{n-1}(P),M)
    \longrightarrow 0 .
\end{split}
\]
\end{corollary}

\begin{proof}
To see the vanishing for \(p>1\), use the length-one flat resolution
of \(H_q(P)\) in the right variable and the chosen projective resolution
of \(M\) in the left variable. The two augmentation comparisons of
Theorem~\ref{res:B05} work for this pair because the terms on each
side are flat. The computation by the right resolution has terms
only in degrees zero and one. Thus all differentials from the second
page onward vanish.
Indeed each differential moves at least two columns.
Only columns \(0,1\) can contain nonzero terms. The finite filtration
on \(H_n\) therefore has \(F_{-1}=0\), \(F_0\) equal to the first
nonzero graded piece, and \(F_1=H_n\). Their identifications are the
terms in the sequence above. The left map can be read at \(F_0\).
A cycle in \(P_n\) tensored with a lift of \(m\in M\) to \(Q_0\)
represents the corresponding class in the filtered total complex,
whose augmentation is the cycle tensored with \(m\). Changing the
lift changes this representative by a vertical boundary, with the
appropriate tensor sign. This gives the usual cycle-tensor edge map.

For example, the length bound holds if every boundary module
\(B_q(P)\) is flat. Exactness of
\(0\to Z_q(P)\to P_q\to B_{q-1}(P)\to0\) and the long exact sequence
of Tor with any left module show \(Z_q(P)\) is flat.
The sequence \(0\to B_q(P)\to Z_q(P)\to H_q(P)\to0\) is then the
required flat resolution. This is the hypothesis and calculation in
\cite[Formula 5.6.5]{weibel}. Additional splitting conclusions use
the hypotheses of the corresponding coefficient theorem, for
example those on the complexes and ring in
\cite[Chapter 11]{p103}.
\end{proof}

\begin{theorem}[Change of rings for Tor]
\label{res:B07}
Let \(f\colon R\to S\) be a homomorphism of rings, \(A\) a right
\(R\)-module and \(M\) a left \(S\)-module. Regard \(M\) as a left
\(R\)-module by restriction. There is a natural first-quadrant
homological spectral sequence
\[
 E^2_{p,q}=\operatorname{Tor}^S_p\bigl(\operatorname{Tor}^R_q(A,S),M\bigr)
       \quad\Longrightarrow\quad \operatorname{Tor}^R_{p+q}(A,M).
\]
In the inner Tor, \(S\) is the \(R\)-\(S\) bimodule obtained from \(f\)
and right multiplication. Its right \(S\)-module structure supplies
the first input of the outer Tor. The target filtration is finite
in each degree.
\end{theorem}

\begin{proof}
Choose a projective right \(R\)-resolution \(P_\bullet\to A\)
and a projective left \(S\)-resolution \(Q_\bullet\to M\).
Use the double complex
\[
    D_{i,j}=P_i\otimes_R Q_j
          \simeq (P_i\otimes_R S)\otimes_S Q_j .
\]
The comparison sends \((x\otimes s)\otimes y\) to \(x\otimes sy\),
with inverse \(x\otimes y\mapsto(x\otimes1)\otimes y\). The forward
map respects the \(R\)-balance relation because
\((xr)\otimes sy=x\otimes f(r)(sy)\), and the \(S\)-balance relation
because \(s(ty)=(st)y\). For the inverse, the corresponding relation is
\[
   ((xr)\otimes1)\otimes y=(x\otimes f(r))\otimes y
                         =(x\otimes 1)\otimes f(r)y.
\]
The two formulas are mutually inverse by \(S\)-balancing and commute
with both differentials.

The augmented complex \(Q_\bullet\to M\) remains exact in left
\(R\)-modules, since restriction of scalars is exact.
For each \(i\) the right module \(P_i\) is \(R\)-flat. Thus the
vertical augmentations \(P_i\otimes_R Q_\bullet\to P_i\otimes_R M\)
are quasi-isomorphisms. Theorem~\ref{res:B04} gives
\[
  H_n(\operatorname{Tot}D)\simeq H_n(P_\bullet\otimes_R M)
          =\operatorname{Tor}^R_n(A,M).
\]
This argument uses the flatness of \(P_i\) over \(R\) and the
exactness of restriction for \(Q_\bullet\). The flatness used for
the other filtration is over \(S\).

Specifically, each \(Q_j\) is \(S\)-flat. With the row index first,
horizontal homology gives
\[
  E^1_{p,q}\simeq H_q(P_\bullet\otimes_R S)\otimes_S Q_p
      =\operatorname{Tor}^R_q(A,S)\otimes_S Q_p .
\]
Here \(P_\bullet\otimes_R S\) is a complex of right \(S\)-modules
whose maps come from those of \(P_\bullet\). Its homology gives
the right action used in the statement. The sign comparison in
Proposition~\ref{res:B03} identifies the \(p\)-direction with the
resolution differential on \(Q\); its homology gives
\(E^2_{p,q}\) by Theorem~\ref{res:B05} over \(S\).
The two indices are nonnegative, so the total degree bounds and the
conclusion of Theorem~\ref{res:B04} apply. For maps in \(A\) and \(M\)
with \(f\) fixed, lifts to their respective resolutions induce maps
of the double complexes. Homotopic choices in \(P\) give the same maps
on the horizontal homology page; homotopic choices in \(Q\) give the
same maps on the homology of that page. Augmentation and the common
total homotopy identities in Theorem~\ref{res:B05} verify naturality
on the target. This proves the change-of-rings comparison of
\cite[Theorem 5.6.6]{weibel}.
\end{proof}
